\begin{trapbox}
\begin{enumerate}
    \item \textbf{不要遗漏斜率不存在的切线}：$y=kx+m$ 不能表示 $x=\pm a$。求\textbf{过定点的所有切线}时，要另查竖直切线，或者直接使用通式 $C^2=a^2A^2+b^2B^2$。

    \item \textbf{相交与相切不可混淆}：$m^2<a^2k^2+b^2$ 表示\textbf{两个不同交点}，不是相切；等号才表示一个交点。对本椭圆，斜率式联立后的二次项系数恒正，所以判别式分类没有退化例外。

    \item \textbf{不要把 $m$ 当成非负数}：截距 $m$ 可以为负。相切时 $m=\pm\sqrt{a^2k^2+b^2}$，所以给定有限斜率总有\textbf{两条}切线；不能只写正号。

    \item \textbf{别把横、纵半轴弄反}：公式的本质是：
    \[
    m^2=(\text{横半轴长})^2k^2+(\text{纵半轴长})^2.
    \]
    如果椭圆写作 $x^2/p^2+y^2/q^2=1$，不论 $p,q$ 哪个更大，都应写 $m^2=p^2k^2+q^2$；不要生搬硬套“$a$ 一定在 $x$ 下”。

    \item \textbf{过定点求斜率可能漏解}：例如椭圆 $x^2/9+y^2/4=1$，从 $P(3,2)$ 引切线，令 $y-2=k(x-3)$ 只能算出 $k=0$，得到 $y=2$。但 $x=3$ 也是经过 $P$ 的切线，必须补上。

    \item \textbf{切点坐标注意符号及前提}：仅当 $y=kx+m$ 已验证相切，才能使用
    \[
    T\left(-\frac{a^2k}{m},\frac{b^2}{m}\right).
    \]
    因为 $a^2k^2+b^2>0$，相切时一定有 $m\ne0$，所以此处不会除以零。

    \item \textbf{椭圆伸缩成圆后，不能混用两个平面的距离}：变换 $x=aX,y=bY$ 只是保持交点个数，并不保留一般的欧氏长度。圆心距公式只能应用于\textbf{变换后的单位圆与直线}。

    \item \textbf{与双曲线公式对比必须补条件}：对于双曲线 $x^2/a^2-y^2/b^2=1$，非竖直切线可满足 $m^2=a^2k^2-b^2$，但还须 $|k|>b/a$。若 $k=\pm b/a,m=0$，直线是\textbf{渐近线而不是切线}，不能把椭圆结论机械照搬过去。

    \item \textbf{切线方程与切点方程不可互相代错}：$\frac{x_0x}{a^2}+\frac{y_0y}{b^2}=1$ 是已知\textbf{切点在椭圆上}时的切线方程，不能把任意椭圆外点直接当作 $(x_0,y_0)$ 来套用。
\end{enumerate}
\end{trapbox}
